\documentclass[UTF8]{ctexart}

% 支持从项目根目录或当前文件目录编译。
\makeatletter
\def\input@path{{../}}
\makeatother
\usepackage{researchnotes}

\title{常见裂项公式}
\author{}
\date{}

\begin{document}
\maketitle

\section{裂项求和的基本思想}

\keyterm{裂项求和}（telescoping summation）的基本思想，是将求和项写成两个表达式之差，
使相邻项或相隔若干项在求和时相互抵消。设 $n$ 为正整数，
$b_1,\ldots,b_{n+1}$ 为实数。若 $a_k=b_k-b_{k+1}$，则
\begin{equation}\label{eq:telescoping-basic}
  \sum_{k=1}^{n}a_k
  =\sum_{k=1}^{n}(b_k-b_{k+1})
  =b_1-b_{n+1}.
\end{equation}
这是因为展开求和后，$b_2,\ldots,b_n$ 各以正、负号出现一次，仅留下首尾两项。
若裂项顺序为 $b_{k+1}-b_k$，结果相应为 $b_{n+1}-b_1$。

下文中的变量与参数均取实数，所有分母均须非零；涉及平方根时，采用实数范围内的非负平方根。
求和上限 $n$ 均为正整数，求和指标 $k$ 取整数。

\section{两个一次因子的乘积}

对于 $h\ne0$，有
\begin{equation}\label{eq:linear-pair}
  \boxed{\frac1{x(x+h)}
  =\frac1h\left(\frac1x-\frac1{x+h}\right)}.
\end{equation}
通分后，括号内的分子为 $(x+h)-x=h$，即得此式。
常见特例为
\begin{align*}
  \frac1{k(k+1)}&=\frac1k-\frac1{k+1},\\
  \frac1{k(k+2)}&=\frac12\left(\frac1k-\frac1{k+2}\right),\\
  \frac1{(2k-1)(2k+1)}
  &=\frac12\left(\frac1{2k-1}-\frac1{2k+1}\right).
\end{align*}
更一般地，对于 $a\ne b$，有
\begin{equation}\label{eq:shifted-linear-pair}
  \boxed{\frac1{(x+a)(x+b)}
  =\frac1{b-a}\left(\frac1{x+a}-\frac1{x+b}\right)}.
\end{equation}

\begin{example}[相邻项抵消]
利用公式 \eqref{eq:linear-pair}，可得
\[
  \sum_{k=1}^{n}\frac1{k(k+1)}
  =\sum_{k=1}^{n}\left(\frac1k-\frac1{k+1}\right)
  =1-\frac1{n+1}=\frac{n}{n+1}.
\]
裂项后，中间的倒数项相互抵消，计算只需保留首尾项。
\end{example}

\section{多个等差因子的乘积}

三个连续因子的裂项公式为
\begin{equation}\label{eq:three-factors}
  \boxed{\frac1{k(k+1)(k+2)}
  =\frac12\left[\frac1{k(k+1)}-\frac1{(k+1)(k+2)}\right]}.
\end{equation}
右侧括号内通分后的分子为 $(k+2)-k=2$。

对于正整数 $m$ 和非零实数 $h$，可推广为
\begin{equation}\label{eq:arithmetic-factors}
  \boxed{
  \frac1{\displaystyle\prod_{j=0}^{m}(x+jh)}
  =\frac1{mh}\left[
  \frac1{\displaystyle\prod_{j=0}^{m-1}(x+jh)}
  -\frac1{\displaystyle\prod_{j=1}^{m}(x+jh)}
  \right]}.
\end{equation}
这里 $\prod$ 表示按所标索引连乘，且要求 $x+jh\ne0$ 对 $j=0,\ldots,m$ 均成立。
将括号内两项通分至左侧分母后，分子为 $(x+mh)-x=mh$，故等式成立。

\begin{example}[三个连续因子的求和]
由公式 \eqref{eq:three-factors}，有
\[
  \sum_{k=1}^{n}\frac1{k(k+1)(k+2)}
  =\frac12\left[\frac1{1\cdot2}-\frac1{(n+1)(n+2)}\right]
  =\frac14-\frac1{2(n+1)(n+2)}.
\]
此时被抵消的是含有两个连续因子的倒数。
\end{example}

\section{平方差型}

由 $x^2-a^2=(x-a)(x+a)$ 及公式 \eqref{eq:shifted-linear-pair}，对于 $a\ne0$，有
\begin{equation}\label{eq:difference-of-squares}
  \boxed{\frac1{x^2-a^2}
  =\frac1{2a}\left(\frac1{x-a}-\frac1{x+a}\right)}.
\end{equation}
常见特例为
\begin{align*}
  \frac1{k^2-1}
  &=\frac12\left(\frac1{k-1}-\frac1{k+1}\right),\qquad k\ge2,\\
  \frac1{4k^2-1}
  &=\frac12\left(\frac1{2k-1}-\frac1{2k+1}\right),\qquad k\ge1.
\end{align*}

\section{分子配成相邻幂之差}

对于正整数 $r$，有
\begin{equation}\label{eq:power-difference}
  \boxed{\frac{(k+1)^r-k^r}{k^r(k+1)^r}
  =\frac1{k^r}-\frac1{(k+1)^r}},\qquad k\ge1.
\end{equation}
该公式由右侧通分直接得到。取 $r=2$ 和 $r=3$，分别得到
\begin{align*}
  \frac{2k+1}{k^2(k+1)^2}
  &=\frac1{k^2}-\frac1{(k+1)^2},\\
  \frac{3k^2+3k+1}{k^3(k+1)^3}
  &=\frac1{k^3}-\frac1{(k+1)^3}.
\end{align*}
因此，对任意正整数 $r$，
\[
  \sum_{k=1}^{n}\frac{(k+1)^r-k^r}{k^r(k+1)^r}
  =1-\frac1{(n+1)^r}.
\]

\section{根式有理化型}

设 $h\ne0$、$x\ge0$ 且 $x+h\ge0$。利用平方差公式，有
\begin{equation}\label{eq:square-root}
  \boxed{\frac1{\sqrt{x+h}+\sqrt{x}}
  =\frac{\sqrt{x+h}-\sqrt{x}}h}.
\end{equation}
上述条件保证分母中的两项不同时为零。取 $x=k$、$h=1$，得到
\[
  \frac1{\sqrt{k+1}+\sqrt{k}}=\sqrt{k+1}-\sqrt{k}.
\]
于是
\[
  \sum_{k=1}^{n}\frac1{\sqrt{k+1}+\sqrt{k}}
  =\sqrt{n+1}-1.
\]

类似地，由立方差公式
\[
  (u-v)(u^2+uv+v^2)=u^3-v^3,
\]
取 $u=\sqrt[3]{k+1}$、$v=\sqrt[3]{k}$，得到
\begin{equation}\label{eq:cube-root}
  \boxed{
  \frac1{\sqrt[3]{(k+1)^2}+\sqrt[3]{k(k+1)}+\sqrt[3]{k^2}}
  =\sqrt[3]{k+1}-\sqrt[3]{k}},\qquad k\ge0.
\end{equation}
这里的立方根均指实立方根，且 $u^3-v^3=1$。

\section{阶乘型}

对于非负整数 $k$，由 $(k+1)!=(k+1)k!$ 及 $0!=1$，有
\begin{equation}\label{eq:factorial}
  \boxed{\frac{k}{(k+1)!}
  =\frac1{k!}-\frac1{(k+1)!}}.
\end{equation}
事实上，右侧通分后的分子为 $(k+1)-1=k$。求和可得
\[
  \sum_{k=1}^{n}\frac{k}{(k+1)!}
  =1-\frac1{(n+1)!}.
\]

\section{三角函数型}

设 $\sin h\ne0$。当 $\sin x\sin(x+h)\ne0$ 时，
\begin{equation}\label{eq:cotangent}
  \boxed{\frac1{\sin x\sin(x+h)}
  =\frac{\cot x-\cot(x+h)}{\sin h}}.
\end{equation}
这是因为
\[
  \cot x-\cot(x+h)
  =\frac{\cos x\sin(x+h)-\sin x\cos(x+h)}{\sin x\sin(x+h)}
  =\frac{\sin h}{\sin x\sin(x+h)}.
\]
当 $\cos x\cos(x+h)\ne0$ 时，同理由正弦差角公式得到
\begin{equation}\label{eq:tangent}
  \boxed{\frac1{\cos x\cos(x+h)}
  =\frac{\tan(x+h)-\tan x}{\sin h}}.
\end{equation}
这两类公式适合角度按等差数列变化的求和。
例如，若 $\sin(x+kh)\ne0$ 对 $k=0,\ldots,n$ 均成立，且 $\sin h\ne0$，则
\[
  \sum_{k=0}^{n-1}\frac1{\sin(x+kh)\sin(x+(k+1)h)}
  =\frac{\cot x-\cot(x+nh)}{\sin h}.
\]

\section{抵消间隔与首尾项的保留}

裂项后应检查两部分的索引间隔。对于正整数 $s$、$n$ 及实数列 $(b_k)_{k\ge1}$，有
\begin{equation}\label{eq:shifted-telescoping}
  \sum_{k=1}^{n}(b_k-b_{k+s})
  =\sum_{k=1}^{s}b_k-\sum_{k=n+1}^{n+s}b_k.
\end{equation}
该式可由两侧分别写成前缀和之差核验。
当 $n\ge s$ 时，抵消后保留首尾各 $s$ 项；当 $n<s$ 时，右侧两个求和区间有重叠，
重叠项仍需相互抵消，公式依然成立。

\begin{example}[间隔为二的裂项]
对正整数 $n$，利用公式 \eqref{eq:linear-pair} 和 \eqref{eq:shifted-telescoping}，得到
\begin{align*}
  \sum_{k=1}^{n}\frac1{k(k+2)}
  &=\frac12\sum_{k=1}^{n}\left(\frac1k-\frac1{k+2}\right)\\
  &=\frac12\left(1+\frac12-\frac1{n+1}-\frac1{n+2}\right).
\end{align*}
当 $n\ge2$ 时，首端保留 $1$、$1/2$，末端保留 $-1/(n+1)$、$-1/(n+2)$。
当 $n=1$ 时，表达式中的两个 $1/2$ 相互抵消，结果同样成立。
\end{example}

应用裂项公式时，可先因式分解或有理化，再确定差项的索引间隔，最后保留未抵消的边界项。
若要求无穷级数的和，还需在有限求和公式的基础上，另行判断边界项是否存在有限极限。

\end{document}
